\subsection{A principal-completion repair of the Nagata example}

\begin{proposition}[Proved; a quasi-excellent principal quotient suffices]
\label{A:principal-completion-bound}
Let $(B,\mathfrak m)$ be a strictly henselian noetherian local ring
of dimension $d$, let $p$ be invertible in $B$, and let
$t\in\mathfrak m$. If $B/tB$ is quasi-excellent, then
\[
 \cd_p\bigl(\Spec B[1/t]\bigr)\leq d.
\]
No quasi-excellence hypothesis on $B$ is imposed.
\end{proposition}
\begin{proof}
Let $C=\varprojlim_n B/t^nB$ be the $t$-adic completion, and let
$g:\Spec C[1/t]\to\Spec B[1/t]$ be its restriction. The pair
$(B,tB)$ is henselian because $B$ is henselian local and
$tB\subset\mathfrak m$. The Fujiwara--Gabber comparison theorem
in the precise unrestricted form proved by
\citet[Theorem 6.11, p.~50 of version 4]{A_BhattMathew2020}
gives, for every torsion abelian \'etale sheaf $\mathcal F$,
\[
 H^q\bigl(\Spec B[1/t],\mathcal F\bigr)
 \simeq
 H^q\bigl(\Spec C[1/t],g^*\mathcal F\bigr)
 \qquad(q\geq0).
\]
In particular this comparison applies to all $p$-primary torsion
coefficients, without a constructibility restriction.

We verify the hypotheses for the affine vanishing theorem on $C$.
The completion is noetherian. The ideal $tC$ is in its Jacobson
radical, and $C/tC=B/tB$ is local, so $C$ is local with maximal
ideal $\mathfrak m C$. A complete ideal-adic pair is henselian;
combining this with the henselian local quotient $B/tB$ shows that
$C$ is henselian. Its residue field equals that of $B$ and is
separably closed. For every $r$, completion of the finite
$B$-module $B/\mathfrak m^r$ along $t$ leaves it unchanged,
because $t$ is nilpotent on this module. Consequently
$C/(\mathfrak m C)^r=B/\mathfrak m^r$ and the maximal-ideal
completion of $C$ is that of $B$. Thus $\dim C=d$.

The lifting theorem of Nishimura--Nishimura and Gabber,
\citet[Theorem 5.1, p.~7]{A_KuranoShimomoto2019}, says that
a noetherian $I$-adically complete ring is quasi-excellent whenever
its quotient by $I$ is quasi-excellent. Applied to $(C,tC)$,
it makes $C$ quasi-excellent. Gabber's strictly local affine
vanishing theorem
\citep[Expos\'e XV, Corollary 1.2.2]{A_ILO2014}
therefore bounds the cohomological dimension of the affine open
$\Spec C[1/t]$ by $d$. The comparison isomorphism gives the
asserted bound on $\Spec B[1/t]$.
\end{proof}

\begin{theorem}[Proved; A.2 for a nonexcellent class of local surfaces]
\label{A:A2-local-surface}
Let $B$ be a strictly henselian noetherian local domain of dimension
two, and let $p$ be invertible in $B$. Suppose that for some
nonzero $t$ in its maximal ideal the ring $B/tB$ is quasi-excellent.
Then A.2 holds for $\Spec B$ and for every nonempty affine open
subscheme of $\Spec B$.

In particular this applies to every regular strictly henselian
two-dimensional local ring of equicharacteristic zero. For the
Nagata--Gabber ring in Theorem~\ref{A:nagata},
every nonempty affine open of $\Spec B[1/(Ft)]$ has
$p$-cohomological dimension two, even though
$\cd_p(\Spec B[1/F])=3$ for the chosen prime of that example.
\end{theorem}
\begin{proof}
Theorem~\ref{A:strict-field} gives
$\cd_p(\Frac B)=2$. Put $R=B[1/t]$.
Proposition~\ref{A:principal-completion-bound} gives
$\cd_p(R)\leq2$. The ring $R$ has dimension at most one:
the only height-two prime of the two-dimensional local domain $B$
is its maximal ideal, which contains $t$. If $R$ is a field,
the conclusion is immediate. Otherwise its normalization
$R^\nu$ is Dedekind by Krull--Akizuki, and exact integral
pushforward gives $\cd_p(R^\nu)\leq2$.
Theorem~\ref{A:Dedekind-Milnor-criterion} shows that the same
upper bound holds on every open of $\Spec R^\nu$.
The normalization comparison in
Theorem~\ref{A:curve-normalization} transfers the upper bound
to every affine open of $\Spec R$.
Here no extra real-field shrinking is needed: if $p=2$,
strict henselianity and $2\in B^\times$ put a square root of
$-1$ in $B$, and hence in every residue field under consideration.
The lower bound two on every nonempty proper open of $\Spec B$
is again Theorem~\ref{A:strict-field}. Thus every nonempty
affine open contained in $\Spec R$ has $p$-cohomological dimension
exactly two.
Every nonempty open of $\Spec B$ meets $\Spec R$ and contains
a nonempty affine open inside that intersection. This is precisely
the quantifier required in A.2, and it also proves the assertion
for any affine open subscheme of $\Spec B$.

If $B$ is regular of equicharacteristic zero, choose $t$ from a
regular system of parameters. Then $B/tB$ is a characteristic-zero
discrete valuation ring, hence excellent
\citep[Expos\'e I, Proposition 3.1]{A_ILO2014}.
This proves the stated class without assuming that $B$ itself
is excellent. In the explicit Nagata--Gabber construction,
the original element $t$ is such a regular parameter.
The nonempty open $\Spec B[1/(Ft)]$ lies in $\Spec B[1/t]$,
so all its nonempty affine subopens have $p$-cohomological dimension
two.
The dimension-three assertion for $\Spec B[1/F]$ was proved
in Theorem~\ref{A:nagata}.
\end{proof}

\begin{remark}[Scope of the comparison]
The completion in this proof is along the principal ideal $(t)$,
not along the maximal ideal. Maximal-ideal completion compares
punctured spectra, whose general upper bound is three in dimension
two, and by itself would not give the needed affine bound.
The proof does not assert that the original nonexcellent surface
becomes quasi-excellent. It proves the cohomological bound on the
specified principal open, and then justifies its inheritance by
smaller opens using the one-dimensional residue criterion.
\end{remark}

\begin{proposition}[Proved; normal models of a quasi-excellent strictly local field]
\label{A:strict-model-rigidity}
Let $B$ be a quasi-excellent strictly henselian noetherian local
domain, let $p$ be invertible in $B$, and put $K=\Frac B$ and
$d=\dim B$. Every noetherian normal domain $R\subset K$ with
$\Frac R=K$ contains $B$ and satisfies
\[
 \cd_p(\Spec R)\leq d=\cd_p(K).
\]
Consequently A.2 holds for all these normal models, in every
dimension. It also holds for every one-dimensional noetherian
domain with fraction field $K$, without a normality assumption.
\end{proposition}
\begin{proof}
We first prove the containment, which is not assumed in the
statement. For every $u\in B^\times$ and integer $r\geq1$,
strict henselianity lifts a simple root of $X^{p^r}-u$ from the
separably closed residue field. Hence $u$ has a $p^r$th root in
$B^\times$. Any discrete valuation $v:K^\times\to\Z$ must
therefore have $v(u)=0$, since this integer is divisible by
every $p^r$. If $b$ belongs to the maximal ideal of $B$, then
$1+b\in B^\times$. The inequality $v(b)<0$ would force
$v(1+b)=v(b)<0$, a contradiction. Thus every element of $B$
belongs to every discrete valuation ring of $K$.

For every height-one prime $\mathfrak q$ of $R$, the ring
$R_{\mathfrak q}$ is a discrete valuation ring. The intersection
formula for a noetherian normal domain
\citep[Tag 031T, Lemma 10.157.6(2)]{A_Stacks} yields
\[
 B\subset\bigcap_{\operatorname{ht}\mathfrak q=1}
 R_{\mathfrak q}=R.
\]
If $R$ is a field, the containment follows from $R=K$.

Write $R$ as the filtered union of its finite-type $B$-subalgebras
$R_i$. Each $R_i$ is a domain, the morphism
$\Spec R_i\to\Spec B$ is affine, finite type and dominant,
and its generic relative dimension is zero, since
$\Frac R_i=K$. The affine bound of
Theorem~\ref{A:relative-affine} gives $\cd_p(R_i)\leq d$.
Descent of constructible sheaves and continuity of their
cohomology, followed by passage to all torsion coefficients as
in Lemma~\ref{A:A2-approximation}, give $\cd_p(R)\leq d$.
The equality for $K$ is Theorem~\ref{A:strict-field}.

Every affine open of $\Spec R$ again has a noetherian normal
coordinate ring with fraction field $K$, so the same upper bound
applies. The generic lower witness from
Proposition~\ref{A:spread} supplies A.2. Finally, for any
one-dimensional noetherian domain with this fraction field,
its normalization is a noetherian Dedekind domain with the same
field; the normal case and Theorem~\ref{A:curve-normalization}
give A.2 without assuming normality.
\end{proof}
